% Generated draft: requires mathematical and source review.
\documentclass[11pt]{article}
\usepackage[margin=25mm]{geometry}
\usepackage{fontspec}
\setmainfont{Latin Modern Roman}
\usepackage{amsmath,amssymb,mathtools,unicode-math}
\setmathfont{Latin Modern Math}
\usepackage{enumitem,needspace,xurl,hyperref}
\setlength{\parindent}{0pt}
\setlength{\parskip}{5pt}
\setlength{\emergencystretch}{3em}
\newcommand{\N}{\mathbb{N}}
\newcommand{\Z}{\mathbb{Z}}
\newcommand{\Q}{\mathbb{Q}}
\newcommand{\R}{\mathbb{R}}
\newcommand{\C}{\mathbb{C}}
\newcommand{\F}{\mathbb{F}}
\newcommand{\A}{\mathbb{A}}
\newcommand{\PP}{\mathbb{P}}
\newcommand{\E}{\mathbb{E}}
\newcommand{\eps}{\varepsilon}
\newcommand{\dd}{\,\mathrm{d}}
\begin{document}
\section*{UnsolvedMath ID: 10900050}
\textbf{Source problem code:} catalogue entry
\subsection{Definitions and mathematical statement}
Retain \(\mathcal O_k\) as in the source's item 7.1, whose statement is not supplied. In the usual number-field interpretation it is the ring of integers of a number field \(k\); any further restriction on \(k\) is unknown. Consider simple undirected graphs on the vertex set \(\mathcal O_k\): distinct \(x,y\) are adjacent in \(\Gamma_{\mathrm{unit}}\) when \(x-y\in\mathcal O_k^\times\), and are adjacent in \(\Gamma_{\mathrm{prime}}\) when their difference is either a unit or a prime. The word ``also'' is retained, so the second graph includes the unit edges. Under the usual prime-element convention, a prime is an element generating a nonzero prime ideal; replacing it by an arbitrary irreducible element could change the graph and requires source clarification.

Write \(\omega(\Gamma)\) for the clique number, the supremum of the cardinalities of vertex subsets whose distinct pairs are adjacent, and put
\[
 q_{\min}(k)=\min_{I\text{ nonzero prime ideal of }\mathcal O_k}|\mathcal O_k/I|
\]
in the finite-residue-field setting intended here. The source records the bound
\[
 \omega(\Gamma_{\mathrm{prime}})
 \leq \omega(\Gamma_{\mathrm{unit}})\,q_{\min}(k).
\]
\emph{Question (Long).} Is this inequality an equality, for \(\mathcal O_k\) in the intended class? That is, does
\[
 \omega(\Gamma_{\mathrm{prime}})
 =\omega(\Gamma_{\mathrm{unit}})\,q_{\min}(k)
\]
hold? The inequality and equality question are one source item. The excerpt does not establish the missing field hypotheses, prime-element convention, finiteness or attainment of clique numbers, or the scope of the quantifier over \(k\). These cannot be resolved by replacing the second graph with a prime-only graph or by changing the recorded bound.
\subsection{Short English statement}
The source gives an upper bound for the clique number of the graph whose differences are units or primes, in terms of the unit graph and smallest prime-ideal residue field. Long asks whether the bound is always sharp in the ring class specified by the missing item 7.1.
\subsection{Sources}
\begin{itemize}
\item Local UnsolvedMath registry, permanent ID 10900050, source code catalogue entry; attribution Long, subsection 7.2, ``Clique numbers in unit- and prime-difference graphs.'' Explicit dependence on item 7.1 is retained.
\item Delp, Hoffoss and Manning, \emph{Problems In Groups, Geometry, and Three-Manifolds} (2015), subsection 7.2, PDF page 6. Supplied URLs: \url{https://arxiv.org/pdf/1512.04620#page=6} and \url{https://arxiv.org/abs/1512.04620}; recorded access 2026-07-29; extraction source-tex checked against rendered PDF text. The inequality and its equality question share the single printed label.
\item Supplied status open and OPEN-TRIAGE. The summary reports no verified result and no specifically addressing publication located in a web search whose cap was reached. This limited search assessment is not an independently verified absence of literature. The 2015 source presents the item as open; no complete modern primary-literature audit is recorded. Status evidence and redistribution terms of the publicly accessible author-submitted TeX and PDF remain NEEDS\_REVIEW.
\end{itemize}
\end{document}
